\documentclass[10pt,a4paper]{amsart}
\usepackage{amsmath,amssymb,amsthm,graphicx}
% Retained preamble for the Stacks Project chapter snapshot in this directory.
% The chapter's own preamble.tex is a copy of the shared Stacks Project preamble; the block
% below is sliced MECHANICALLY, line for line, from it. Nothing was added, removed or
% redefined; the wrapper class, the xr-hyper/hyperref wiring and the externaldocument list
% are omitted because the assembler supplies the document class.
%
% lines taken from preamble.tex: 3, 5, 9, 11, 12, 156, 157, 158, ... (51 lines total)
%
% package coverage restored from the chapter's own preamble (2026-10-13): xy, xy-2cell

\usepackage[all]{xy}
\xyoption{2cell}
\UseAllTwocells

\newenvironment{reference}{\comment}{\endcomment}
\newenvironment{slogan}{\comment}{\endcomment}
\newenvironment{history}{\comment}{\endcomment}
\theoremstyle{plain}
\newtheorem{theorem}[subsection]{Theorem}
\newtheorem{proposition}[subsection]{Proposition}
\newtheorem{lemma}[subsection]{Lemma}
\theoremstyle{definition}
\newtheorem{definition}[subsection]{Definition}
\newtheorem{example}[subsection]{Example}
\newtheorem{exercise}[subsection]{Exercise}
\newtheorem{situation}[subsection]{Situation}
\theoremstyle{remark}
\newtheorem{remark}[subsection]{Remark}
\newtheorem{remarks}[subsection]{Remarks}
\numberwithin{equation}{subsection}
\def\lim{\mathop{\mathrm{lim}}\nolimits}
\def\colim{\mathop{\mathrm{colim}}\nolimits}
\def\Spec{\mathop{\mathrm{Spec}}}
\def\Hom{\mathop{\mathrm{Hom}}\nolimits}
\def\Ext{\mathop{\mathrm{Ext}}\nolimits}
\def\SheafHom{\mathop{\mathcal{H}\!\mathit{om}}\nolimits}
\def\SheafExt{\mathop{\mathcal{E}\!\mathit{xt}}\nolimits}
\def\Sch{\mathit{Sch}}
\def\Mor{\mathop{\mathrm{Mor}}\nolimits}
\def\Ob{\mathop{\mathrm{Ob}}\nolimits}
\def\Sh{\mathop{\mathit{Sh}}\nolimits}
\def\NL{\mathop{N\!L}\nolimits}
\def\CH{\mathop{\mathrm{CH}}\nolimits}
\def\proetale{{pro\text{-}\acute{e}tale}}
\def\etale{{\acute{e}tale}}
\def\QCoh{\mathit{QCoh}}
\def\Ker{\mathop{\mathrm{Ker}}}
\def\Im{\mathop{\mathrm{Im}}}
\def\Coker{\mathop{\mathrm{Coker}}}
\def\Coim{\mathop{\mathrm{Coim}}}
\DeclareMathSymbol{\boxtimes}{\mathbin}{AMSa}{"02}
\def\QCohstack{\mathcal{QC}\!\mathit{oh}}
\def\Cohstack{\mathcal{C}\!\mathit{oh}}
\def\Spacesstack{\mathcal{S}\!\mathit{paces}}
\def\Quotfunctor{\mathrm{Quot}}
\def\Hilbfunctor{\mathrm{Hilb}}
\def\Curvesstack{\mathcal{C}\!\mathit{urves}}
\def\Polarizedstack{\mathcal{P}\!\mathit{olarized}}
\def\Complexesstack{\mathcal{C}\!\mathit{omplexes}}
\def\Pic{\mathop{\mathrm{Pic}}\nolimits}
\def\Picardstack{\mathcal{P}\!\mathit{ic}}
\def\Picardfunctor{\mathrm{Pic}}
\def\Deformationcategory{\mathcal{D}\!\mathit{ef}}
\usepackage{cleveref}
\crefname{originallemma}{original lemma}{original lemmas}
\Crefname{originallemma}{Original lemma}{Original lemmas}
\makeatletter
\newlabel{formal-defos-lemma-length}{{06GG}{source}{}{sourcepointer-1}{}}
\newlabel{formal-defos-lemma-length@cref}{{[originallemma][0][]06GG}{[1][0][]source}{}{}{}}
\makeatother
\crefname{originallemma}{original lemma}{original lemmas}
\Crefname{originallemma}{Original lemma}{Original lemmas}
\makeatletter
\newlabel{formal-defos-lemma-surjective-cotangent-space}{{06GZ}{source}{}{sourcepointer-2}{}}
\newlabel{formal-defos-lemma-surjective-cotangent-space@cref}{{[originallemma][0][]06GZ}{[1][0][]source}{}{}{}}
\makeatother
\crefname{originallemma}{original lemma}{original lemmas}
\Crefname{originallemma}{Original lemma}{Original lemmas}
\makeatletter
\newlabel{formal-defos-remark-versal-object}{{06HS}{source}{}{sourcepointer-3}{}}
\newlabel{formal-defos-remark-versal-object@cref}{{[originallemma][0][]06HS}{[1][0][]source}{}{}{}}
\makeatother
\crefname{originallemma}{original lemma}{original lemmas}
\Crefname{originallemma}{Original lemma}{Original lemmas}
\makeatletter
\newlabel{lemma-approximate}{{07XB}{source}{}{sourcepointer-4}{}}
\newlabel{lemma-approximate@cref}{{[originallemma][0][]07XB}{[1][0][]source}{}{}{}}
\makeatother
\crefname{originallemma}{original lemma}{original lemmas}
\Crefname{originallemma}{Original lemma}{Original lemmas}
\makeatletter
\newlabel{lemma-versality-matches}{{0CXK}{source}{}{sourcepointer-5}{}}
\newlabel{lemma-versality-matches@cref}{{[originallemma][0][]0CXK}{[1][0][]source}{}{}{}}
\makeatother
\title{Producing a versal finite-type family from a versal formal object}
\author{The Stacks Project Authors}
\date{stacks-project:artin.tex:local-snapshot-20261012; GNU FDL 1.2 or later}
\begin{document}
\maketitle
\noindent\textit{Editorial note.} Statement span is lines 2091-2107 of sources/stacks-artin/pinned-original-source.tex; proof span is lines 2109-2149, the complete author proof (39 body lines).\par\medskip
\noindent\textit{Editorial note.} Navigation stubs: 5 cited-but-not-included labels are supplied as non-load-bearing pointers with their pinned tags where the snapshot tag index has them.\par\medskip
\noindent\textit{Editorial note.} extra\_packages is empty on purpose: the retained preamble already loads xy with options and a second declaration clashes in the assembler wrapper.\par\medskip
\noindent\textit{Editorial note.} Text is retained verbatim, including the author's own typos and placeholders.\par\medskip
\noindent\textit{Editorial note.} Queue index 26 of queue\_artin.json (pb=1616, stmt 2091-2107, proof 2109-2149).\par\medskip
\noindent\textit{Editorial note.} stacks-artin wave ar01 (2026-10-14): fresh adjudication accept per lane-s3/artin\_judgments.md; pool per Chief 10:5x direct dispatch (idle-first over S12's 01:5x note, SMC precedent); artin.tex 99 rows/67 with proof judged 46 accept / 21 reject / 5 AGP-overlap avoidance rows; patch-sd-envs contract spans scanned.\par\medskip
\noindent\textit{Allowed standard prerequisites (editorial).}\begin{itemize}
\item Explicit prerequisite boundary. The cited inputs are prerequisites; the argument in the retained span is not replaced by any cited result.
\end{itemize}\medskip
\ifdefined\hypertarget\hypertarget{sourcepointer-1}{}\fi
\noindent\textit{Source reference.} Original lemma 06GG: sources/stacks-artin/pinned-original-source.tex, source lines 2149--2149. Cited by the retained statement or proof but not part of the collected spans; supplied only as a navigation pointer. The argument does not depend on its content.\par\medskip
\ifdefined\hypertarget\hypertarget{sourcepointer-2}{}\fi
\noindent\textit{Source reference.} Original lemma 06GZ: sources/stacks-artin/pinned-original-source.tex, source lines 2149--2149. Cited by the retained statement or proof but not part of the collected spans; supplied only as a navigation pointer. The argument does not depend on its content.\par\medskip
\ifdefined\hypertarget\hypertarget{sourcepointer-3}{}\fi
\noindent\textit{Source reference.} Original lemma 06HS: sources/stacks-artin/pinned-original-source.tex, source lines 2149--2149. Cited by the retained statement or proof but not part of the collected spans; supplied only as a navigation pointer. The argument does not depend on its content.\par\medskip
\ifdefined\hypertarget\hypertarget{sourcepointer-4}{}\fi
\noindent\textit{Source reference.} Original lemma 07XB: sources/stacks-artin/pinned-original-source.tex, source lines 1488--1508. Cited by the retained statement or proof but not part of the collected spans; supplied only as a navigation pointer. The argument does not depend on its content.\par\medskip
\ifdefined\hypertarget\hypertarget{sourcepointer-5}{}\fi
\noindent\textit{Source reference.} Original lemma 0CXK: sources/stacks-artin/pinned-original-source.tex, source lines 1903--1914. Cited by the retained statement or proof but not part of the collected spans; supplied only as a navigation pointer. The argument does not depend on its content.\par\medskip
\begin{lemma}
\label{lemma-approximate-versal}
Let $S$ be a locally Noetherian scheme. Let
$p : \mathcal{X} \to (\Sch/S)_{fppf}$ be a category fibred in groupoids.
Let $\xi = (R, \xi_n, f_n)$ be a formal object of $\mathcal{X}$ with
$\xi_1$ lying over $\Spec(k) \to S$ with image $s \in S$. Assume
\begin{enumerate}
\item $\xi$ is versal,
\item $\xi$ is effective,
\item $\mathcal{O}_{S, s}$ is a G-ring, and
\item $p : \mathcal{X} \to (\Sch/S)_{fppf}$ is limit preserving on objects.
\end{enumerate}
Then there exist a morphism of finite type $U \to S$, a finite type
point $u_0 \in U$ with residue field $k$, and an object $x$ of $\mathcal{X}$
over $U$ such that $x$ is versal at $u_0$ and such that
$x|_{\Spec(\mathcal{O}_{U, u_0}/\mathfrak m_{u_0}^n)} \cong \xi_n$.
\end{lemma}

\begin{proof}
Choose an object $x_R$ of $\mathcal{X}$ lying over $\Spec(R)$ whose associated
formal object is $\xi$. Let $N = 2$ and apply Lemma \ref{lemma-approximate}.
We obtain $A, \mathfrak m_A, x_A, \ldots$.
Let $\eta = (A^\wedge, \eta_n, g_n)$ be the formal object associated to
$x_A|_{\Spec(A^\wedge)}$. We have a diagram
$$
\vcenter{
\xymatrix{
& \eta \ar[d] \\
\xi \ar[r] \ar@{..>}[ru] & \xi_2 = \eta_2
}
}
\quad\text{lying over}\quad
\vcenter{
\xymatrix{
& A^\wedge \ar[d] \\
R \ar[r] \ar@{..>}[ru] & R/\mathfrak m_R^2 = A/\mathfrak m_A^2
}
}
$$
The versality of $\xi$ means exactly that we can find the
dotted arrows in the diagrams, because we can successively find
morphisms $\xi \to \eta_3$, $\xi \to \eta_4$, and so on by
Formal Deformation Theory, Remark \ref{formal-defos-remark-versal-object}.
The corresponding ring map $R \to A^\wedge$ is surjective by
Formal Deformation Theory, Lemma
\ref{formal-defos-lemma-surjective-cotangent-space}.
On the other hand, we have
$\dim_k \mathfrak m_R^n/\mathfrak m_R^{n + 1} =
\dim_k \mathfrak m_A^n/\mathfrak m_A^{n + 1}$ for all $n$ by construction.
Hence $R/\mathfrak m_R^n$ and $A/\mathfrak m_A^n$ have the same (finite)
length as $\Lambda$-modules by additivity of length and
Formal Deformation Theory, Lemma \ref{formal-defos-lemma-length}.
It follows that $R/\mathfrak m_R^n \to A/\mathfrak m_A^n$ is an isomorphism
for all $n$, hence $R \to A^\wedge$ is an isomorphism. Thus $\eta$ is
isomorphic to a versal object, hence versal itself. By
Lemma \ref{lemma-versality-matches}
we conclude that $x_A$ is versal at the point $u_0$ of
$U = \Spec(A)$ corresponding to $\mathfrak m_A$.
\end{proof}

\end{document}
